\documentclass[10pt,a4paper]{amsart}
\usepackage[margin=19mm]{geometry}
\usepackage{amsmath,amssymb,mathtools}
\usepackage[scr=rsfs,cal=euler]{mathalfa}
\usepackage{xfrac}
\usepackage[unicode,hidelinks,bookmarks=false]{hyperref}
\newcommand{\ie}{i.e.\ }
\newcommand{\cf}{cf.\ }
\newcommand{\resp}{respectively\ }
\newcommand{\Subsection}{\subsection}
\providecommand{\nobreakdash}{\nobreak\hbox{-}}
\setlength{\emergencystretch}{2em}
\multlinegap0pt
\usepackage{amssymb}
\usepackage{xcolor}
\usepackage[shortlabels]{enumitem}
\usepackage{tikz}
\usepackage{natbib}
\newtheorem{theorem}{Theorem}
\newtheorem{proposition}{Proposition}
\def\top{{ \mathrm{\scriptscriptstyle T} }}
\title{Identifiability of Treatment Effects with Unobserved Spatially Varying Confounders}
\author{Tommy Tang, Xinran Li, Bo Li}
\date{Source: arXiv:2602.23291v1 (CC BY 4.0)}
\begin{document}
\maketitle

\noindent\emph{Source and licence.} Identifiability of Treatment Effects with Unobserved Spatially Varying Confounders. Version arXiv:2602.23291v1 (CC BY 4.0), distributed by arXiv under the Creative Commons Attribution 4.0 International licence, http://creativecommons.org/licenses/by/4.0/, as recorded in the arXiv licence metadata for this exact version and shown on the abstract page https://arxiv.org/abs/2602.23291v1. Full licence notice: \texttt{licenses/arxiv-2602.23291-CC-BY-4.0.txt}.\par\smallskip

\noindent\textit{Editorial scope.} Counted unit: the article's proposition prop:car-fullyconnected-nonid (source0.tex lines 1451--1459). The article's complete original proof of it is delivered with it (source0.tex lines 1461--1625, one \texttt{proof} environment of 165 lines). No mathematics is written, repaired or re-derived: the statement and the proof are the article's own lines, in the article's own notation, and no retained line is substituted or re-flowed.  Retained uncounted as the source-local context the proof needs: lines 171--190; lines 228--234; lines 257--266; lines 1193--1209; lines 1371--1413.  Nothing was omitted from any retained span, so the package carries no omission record.  No retained line contains a citation, so the wrapper carries no bibliography and no bibliography entry.  No retained line contains a figure environment, an image-inclusion command or drawing code, so no image is delivered and the package declares no asset dependency.  Editorial formatting only, in addition: an AMS \texttt{amsart} preamble that restates the article's own declarations and macros used by the retained lines, namely the environments \texttt{theorem}, \texttt{proposition} on separate counters, together with the standard packages those lines need.  The retained environments print in this document as Theorem 1, Proposition 1 in order of appearance, whereas the article prints the same items with the numbering of its own full document.  The complete native source of the article as distributed in the arXiv e-print for this exact version is bundled unchanged as \texttt{sources/195390/source0.tex}.
\begin{equation}\label{eq:simp-model}
    Y = Z \beta + U + \epsilon, 
    \quad \epsilon \sim \mathcal{N}(0, \sigma_{\epsilon}^2 I_n), 
    \quad 
    \begin{pmatrix}
        U\\
        Z
    \end{pmatrix} 
    \sim \mathcal{N}
    \left(  
    \begin{pmatrix}
        0\\
        0
    \end{pmatrix}, 
    \begin{pmatrix}
        \Sigma_{UU} & \Sigma_{UZ}\\
        \Sigma_{ZU} & \Sigma_{ZZ}
    \end{pmatrix}
    \right). 
\end{equation}

\begin{align}\label{eq:Sigma_SP}
    \Sigma = 
    \begin{pmatrix} 
        \tau_U(D-\varphi_U W) & -\rho\sqrt{\tau_U\tau_Z} D \\
        -\rho\sqrt{\tau_U\tau_Z} D & \tau_Z(D-\varphi_Z W)
    \end{pmatrix}^{-1}.
\end{align}

\begin{theorem}\label{thm:car-identifiability}
Consider $n$ spatial locations with a proximity matrix $W$. 
Let 
$\mathcal{M}_1,...,\mathcal{M}_B$ be the disjoint connected components, which form a partition of all the locations. 
Assume that the data generating process follows \eqref{eq:simp-model} and \eqref{eq:Sigma_SP}. 
If $\varphi_U
\neq
0$ and there exists one component $\mathcal{M}_b$ such that 
either the corresponding proximity matrix $W_{[b]}$ has at least three distinct eigenvalues or the corresponding $D_{[b]}$ contains distinct diagonal elements, then all the model parameters $\tau_U, \tau_Z, \varphi_U, \varphi_Z, \rho, \sigma_\epsilon^2$ and $\beta$ are identifiable. 
\end{theorem}

\begin{align}
    \beta I + \rho \sqrt{\tau_Z/\tau_U} (I-\varphi_U \Lambda)^{-1} & = 
    \tilde{\beta} I + \tilde{\rho} \sqrt{\tilde{\tau}_Z/\tilde{\tau}_U} (I-\tilde{\varphi}_U \Lambda)^{-1}, 
    \label{eq:iden_cond_mean_Y_Z}
    \\ 
    \tau_U^{-1} (I-\varphi_U \Lambda)^{-1} 
    +\tau_{\epsilon}^{-1} \Gamma^\top D \Gamma
    & = 
    \tilde{\tau}_U^{-1} (I-\tilde{\varphi}_U \Lambda)^{-1} 
    +\tilde{\tau}_{\epsilon}^{-1} \Gamma^\top D \Gamma, 
    \label{eq:iden_var_mean_Y_Z}
    \\ 
    \tau_Z \big\{ I - \varphi_Z \Lambda - \rho^2 (I - \varphi_U \Lambda)^{-1} \big\}
    & = 
    \tilde{\tau}_Z \big\{ I - \tilde{\varphi}_Z \Lambda - \tilde{\rho}^2 (I - \tilde{\varphi}_U \Lambda)^{-1} \big\}.
    \label{eq:prec_Z}
\end{align}

\begin{align*}
    & D^{-1/2} \Gamma
    \big\{ \beta I +\rho \sqrt{\tau_Z/\tau_U} (I-\varphi_U \Lambda)^{-1} \big\} 
    \Gamma^{\top} D^{1/2}
    = aI \text{ for some } a
    \\
    \Longleftrightarrow \quad  
    & 
    \beta I +\rho \sqrt{\tau_Z/\tau_U} (I-\varphi_U \Lambda)^{-1} 
    = aI \text{ for some } a
    \\
    \Longleftrightarrow \quad  
    & 
    \beta I +\rho \sqrt{\tau_Z/\tau_U} (I-\varphi_U \Lambda)^{-1} 
    = aI \text{ for some } a
    \\
    \Longleftrightarrow \quad  
    & 
    I-\varphi_U \Lambda
    = aI \text{ for some } a, \text{ or } \rho = 0\\
    \Longleftrightarrow \quad  
    & 
    \varphi_U = 0 \text{ or } \rho = 0. 
\end{align*}


\subsection{A remark on the necessity of $\varphi_U\neq 0$}\label{sec:nece_phiU_car}
Let $(\beta, \tau_Z, \varphi_Z, \tau_U, \varphi_U, \tau_{\epsilon}, \rho)$ be the true data-generating parameters, where we again use $\tau_{\epsilon}$ to denote $\sigma_{\epsilon}^{-2}$ for convenience. 
Below we show that when $\varphi_U = 0$, the model parameters are not identifiable. 
From the proof of Theorem \ref{thm:car-identifiability}, 
to prove nonidentifiability, 
it suffices to find parameters 
$(\tilde{\beta}, \tilde{\tau}_Z, \tilde{\varphi}_Z, \tilde{\tau}_U, \tilde{\varphi}_U, \tilde{\tau}_{\epsilon}, \tilde{\rho})$
that are different from the true data-generating parameters and yield the same values as in \eqref{eq:iden_cond_mean_Y_Z}--\eqref{eq:prec_Z}. 
Let $\tilde{\varphi}_U=\varphi_U=0$, $\tilde{\tau}_U = \tau_U$ and $\tilde{\tau}_\epsilon = \tau_\epsilon$. Then \eqref{eq:iden_cond_mean_Y_Z}--\eqref{eq:prec_Z} hold if and only if 
\begin{align*}
    \beta I + \rho \sqrt{\tau_Z/\tau_U} I & = 
    \tilde{\beta} I + \tilde{\rho} \sqrt{\tilde{\tau}_Z/\tau_U} I, 
    \\ 
    \tau_Z \big( I - \varphi_Z \Lambda - \rho^2 I \big)
    & = 
    \tilde{\tau}_Z \big( I - \tilde{\varphi}_Z \Lambda - \tilde{\rho}^2 I \big).
\end{align*}

\begin{proposition}\label{prop:car-fullyconnected-nonid}
Consider the same model setup as in Theorem \ref{thm:car-identifiability}. 
If 
$\rho \ne 0$, and
$W$ corresponds to a fully connected neighborhood structure, i.e., the diagonal elements of $W$ are all $0$ and the off-diagonal elements of $W$ are all $1$, 
then the treatment effect $\beta$ a
is 
not identifiable.
\end{proposition}

\begin{proof}[Proof of Proposition \ref{prop:car-fullyconnected-nonid}]
Suppose that $(\beta, \varphi_U, \tau_U, \varphi_Z, \tau_Z, \rho, \tau_\epsilon)$ denote the true data generating parameter, where we again use $\tau_{\epsilon}$ to denote $\sigma_{\epsilon}^{-2}$ for convenience. 
When $\varphi_U = 0$, Proposition \ref{prop:car-fullyconnected-nonid} follows immediately from Section \ref{sec:nece_phiU_car}. 
Below we consider only the case where $\varphi_U \ne 0$. 
From the proof of Theorem \ref{thm:car-identifiability}, it suffices to find other parameters $(\tilde\beta, \tilde\varphi_U, \tilde\tau_U, \tilde\varphi_Z, \tilde\tau_Z, \tilde\rho, \tilde\tau_\epsilon)$ such that equations \eqref{eq:iden_cond_mean_Y_Z}--\eqref{eq:prec_Z} hold and $\tilde\beta \ne \beta$. 

First, we simplify \eqref{eq:iden_cond_mean_Y_Z}--\eqref{eq:prec_Z} using the condition that $W$ represents a connected neighborhood structure. 
By definition, we have $W = 1_{n} 1_{n}^\top - I_n$ and $D = (n-1)I_n$. 
Because the eigenvalues of $1_{n} 1_{n}^\top$ are either $n$ or $0$, 
the eigenvalues of $W$ must be either $n-1$ or $-1$. 
This further implies that the eigenvalue of $D^{-1/2}WD^{-1/2} = (n-1)^{-1} W$ are either $1$ or $-(n-1)^{-1}$. 
We can then verify that equations \eqref{eq:iden_cond_mean_Y_Z}--\eqref{eq:prec_Z}  are equivalent to 

\begin{align}
\beta + \rho\sqrt{\tau_Z/\tau_U}(1-\varphi_U)^{-1} &= \tilde\beta + \tilde\rho\sqrt{\tilde\tau_Z/\tilde\tau_U}(1-\tilde\varphi_U)^{-1}\label{eq:fullyconnected-exp-product-1} \\
\beta + \rho\sqrt{\tau_Z/\tau_U}\left(1+\frac{\varphi_U}{n-1}\right)^{-1} &= \tilde\beta + \tilde\rho\sqrt{\tilde\tau_Z/\tilde\tau_U}\left(1+\frac{\tilde\varphi_U}{n-1}\right)^{-1}\label{eq:fullyconnected-exp-product-n-1} 
\end{align}
\begin{align}
\tau_U^{-1}(1-\varphi_U)^{-1} + \tau_\epsilon^{-1} (n-1) &= \tilde\tau_U^{-1}(1-\tilde\varphi_U)^{-1} + \tilde\tau_\epsilon^{-1}(n-1)\label{eq:fullyconnected-condvar-1} \\
\tau_U^{-1}\left(1+\frac{\varphi_U}{n-1}\right)^{-1}  + \tau_\epsilon^{-1} (n-1) &= \tilde\tau_U^{-1}\left(1+\frac{\tilde\varphi_U}{n-1}\right)^{-1}  + \tilde\tau_\epsilon^{-1}(n-1)\label{eq:fullyconnected-condvar-n-1} 
\end{align}
\begin{align}
\tau_Z\left\{1 - \varphi_Z - \rho^2(1-\varphi_U)^{-1}\right\} &= \tilde\tau_Z\left\{1 - \tilde\varphi_Z - \tilde\rho^2(1-\tilde\varphi_U)^{-1}\right\}\label{eq:fullyconnected-precision-1} \\
\tau_Z\left\{1 + \frac{\varphi_Z}{n-1} - \rho^2\left(1+\frac{\varphi_U}{n-1}\right)^{-1}\right\} &= \tilde\tau_Z\left\{1 + \frac{\tilde\varphi_Z}{n-1} - \tilde\rho^2\left(1+\frac{\tilde\varphi_U}{n-1}\right)^{-1}\right\}\label{eq:fullyconnected-precision-n-1} 
\end{align}



Second, we try to construct parameters $(\tilde\beta, \tilde\varphi_U, \tilde\tau_U, \tilde\varphi_Z, \tilde\tau_Z, \tilde\rho, \tilde\tau_\epsilon)$ such that equations \eqref{eq:fullyconnected-exp-product-1}--\eqref{eq:fullyconnected-precision-n-1} hold. 
Let $b$ be a scalar whose value will be determined later. Define 
\begin{align}
r(b) &= \frac{\left(1 + \frac{\varphi_U}{n-1} \right)^{-1} - \left(1-\varphi_U\right)^{-1}}{\left(1 + \frac{b\varphi_U}{n-1} \right)^{-1} - \left(1-b\varphi_U\right)^{-1}} 
,
\label{eq:car-full-r}
\\
\nonumber
B_1(b) & = \tau_Z + \frac{\tau_Z\varphi_Z}{n-1} - \rho^2\tau_Z\left( 1 + \frac{\varphi_U}{n-1}\right)^{-1} + r\rho^2 \tau_Z \left( 1 + \frac{b\varphi_U}{n-1}\right)^{-1},
\\
\nonumber
B_2(b) & = \tau_Z - \tau_Z \varphi_Z - \rho^2\tau_Z(1-\varphi_U)^{-1} + r\rho^2\tau_Z(1-b\varphi_U)^{-1}.
\end{align}
For notational convenience, we will simply write $r(b), B_1(b)$ and $B_2(b)$ as  $r, B_1$ and $B_2$. 
Define further $(\tilde\beta, \tilde\varphi_U, \tilde\tau_U, \tilde\varphi_Z, \tilde\tau_Z, \tilde\rho, \tilde\tau_\epsilon)$ as follows:
\begin{align}\label{eq:tilde_car_non}
    \tilde\varphi_U & = b\varphi_U, 
    \qquad 
    \tilde\tau_U^{-1} = r \tau_U^{-1}
    \qquad 
    \tilde\tau_\epsilon^{-1} = (n-1)^{-1}\tau_U^{-1}\{(1-\varphi_U)^{-1} - r(1-b\varphi_U)^{-1}\} + \tau_\epsilon^{-1}, 
    \nonumber
    \\
    \tilde\beta & = \beta + \rho\sqrt{\frac{\tau_Z}{\tau_U}}(1-\varphi_U)^{-1} - r\rho\sqrt{\frac{\tau_Z}{\tau_U}}(1-b\varphi_U)^{-1}, 
    \qquad 
    \tilde\varphi_Z = \frac{B_1 - B_2}{(n-1)^{-1}B_2 + B_1},
    \nonumber
    \\
    \tilde\tau_Z &= 
    \frac{B_2}{1-\tilde\varphi_Z}
    = 
    \frac{(n-1)B_1 + B_2}{n}, \qquad \tilde\rho = \rho\sqrt{r \tau_Z/\tilde\tau_Z}. 
\end{align}

Third, we show that equations \eqref{eq:fullyconnected-exp-product-1}--\eqref{eq:fullyconnected-precision-n-1} hold for $(\tilde\beta, \tilde\varphi_U, \tilde\tau_U, \tilde\varphi_Z, \tilde\tau_Z, \tilde\rho, \tilde\tau_\epsilon)$ defined in \eqref{eq:tilde_car_non}, under the assumption that the quantities in \eqref{eq:tilde_car_non} are well-defined. 
We begin with \eqref{eq:fullyconnected-condvar-1} and \eqref{eq:fullyconnected-condvar-n-1}. 
By the definition in \eqref{eq:car-full-r} and \eqref{eq:tilde_car_non}, we have 
\begin{align*}
\tilde\tau_U^{-1}(1-\tilde\varphi_U)^{-1} + \tilde\tau_\epsilon^{-1}(n-1) &= r \tau_U^{-1}(1-b\varphi_U)^{-1} + \tau_U^{-1}\{(1-\varphi_U)^{-1} - r(1-b\varphi_U)^{-1}\} + \tau_\epsilon^{-1}(n-1)
\\
&= \tau_U^{-1}(1-\varphi_U)^{-1} + \tau_\epsilon^{-1}(n-1)
\end{align*} and 
\begin{align*}
& \quad \ \tilde\tau_U^{-1}\left(1+\frac{\tilde\varphi_U}{n-1}\right)^{-1} + \tilde\tau_\epsilon^{-1}(n-1)
\\ 
&= r \tau_U^{-1}\left(1+\frac{b\varphi_U}{n-1}\right)^{-1} + \tau_U^{-1}\{(1-\varphi_U)^{-1} - r(1-b\varphi_U)^{-1}\} + \tau_\epsilon^{-1}(n-1) \\
&= r\tau_U^{-1}\left\{\left(1+\frac{b\varphi_U}{n-1}\right)^{-1} - (1-b\varphi_U)^{-1}\right\} + \tau_U^{-1}(1-\varphi_U)^{-1} + \tau_\epsilon^{-1}(n-1) \\
&= \tau_U^{-1}\left\{\left(1 + \frac{\varphi_U}{n-1} \right)^{-1} - \left(1-\varphi_U\right)^{-1}\right\}+ \tau_U^{-1}(1-\varphi_U)^{-1} + \tau_\epsilon^{-1}(n-1) \\
&= \tau_U^{-1}\left(1+\frac{\varphi_U}{n-1}\right)^{-1} + \tau_\epsilon^{-1} (n-1),
\end{align*} 
which imply that \eqref{eq:fullyconnected-condvar-1} and \eqref{eq:fullyconnected-condvar-n-1} hold. We next consider \eqref{eq:fullyconnected-exp-product-1} and \eqref{eq:fullyconnected-exp-product-n-1}. 
By definition, we have  \begin{equation}\nonumber%
\tilde\rho\sqrt{\tilde\tau_Z/\tilde\tau_U}=
 \rho\sqrt{r \tau_Z/\tilde\tau_Z} \sqrt{\tilde\tau_Z/\tilde\tau_U}
=
\rho\sqrt{r \tau_Z/\tilde\tau_U}
= 
r\rho\sqrt{\tau_Z/\tau_U}.
\end{equation} 
Consequently, 
\begin{align*}
&\tilde\beta + \tilde\rho\sqrt{\tilde\tau_Z/\tilde\tau_U}(1-\tilde\varphi_U)^{-1}\\ &=  \beta + \rho\sqrt{\frac{\tau_Z}{\tau_U}}(1-\varphi_U)^{-1} - r\rho\sqrt{\frac{\tau_Z}{\tau_U}}(1-b\varphi_U)^{-1} + r\rho\sqrt{\tau_Z/\tau_U}(1-b\varphi_U)^{-1} \\
&= \beta + \rho\sqrt{\tau_Z/\tau_U}(1-\varphi_U)^{-1},
\end{align*} 
and 
\begin{align*}
& \quad \ \tilde\beta + \tilde\rho\sqrt{\tilde\tau_Z/\tilde\tau_U}\left(1+\frac{\tilde\varphi_U}{n-1}\right)^{-1} 
\\
&= \beta + \rho\sqrt{\frac{\tau_Z}{\tau_U}}(1-\varphi_U)^{-1} - r\rho\sqrt{\frac{\tau_Z}{\tau_U}}(1-b\varphi_U)^{-1} + r\rho\sqrt{\frac{\tau_Z}{\tau_U}}\left(1+ \frac{b\varphi_U}{n-1} \right)^{-1} \\
&= \beta + \rho\sqrt{\frac{\tau_Z}{\tau_U}}(1-\varphi_U)^{-1} + r\rho \sqrt{\frac{\tau_Z}{\tau_U}}\left\{\left(1+ \frac{b\varphi_U}{n-1} \right)^{-1}-(1-b\varphi_U)^{-1} \right\} \\
&= \beta + \rho\sqrt{\frac{\tau_Z}{\tau_U}}(1-\varphi_U)^{-1} + \rho\sqrt{\frac{\tau_Z}{\tau_U}}\left\{\left(1+ \frac{\varphi_U}{n-1} \right)^{-1}-(1-\varphi_U)^{-1}\right\}\\
&= \beta +\rho\sqrt{\frac{\tau_Z}{\tau_U}}\left(1+ \frac{\varphi_U}{n-1} \right)^{-1},
\end{align*} 
which imply that \eqref{eq:fullyconnected-exp-product-1} and \eqref{eq:fullyconnected-exp-product-n-1} hold. 
We finally consider \eqref{eq:fullyconnected-precision-1}
and \eqref{eq:fullyconnected-precision-n-1}. 
By definition, 
\begin{align*}
\tilde\tau_Z\{1-\tilde\varphi_Z - \tilde\rho^2(1-\tilde\tau_U)^{-1}\} &= \frac{B_2}{1-\tilde\varphi_Z}(1-\tilde\varphi_Z) - \tilde\tau_Z\tilde\rho^2(1-b\varphi_U)^{-1}
\\ &= B_2 - r \tau_Z \rho^2 (1-b\varphi_U)^{-1} \\
&= \tau_Z - \tau_Z \varphi_Z - \rho^2\tau_Z(1-\varphi_U)^{-1} 
\\ &= \tau_Z \{ 1 - \varphi_Z - \rho^2(1-\varphi_U)^{-1} \}, 
\end{align*} 
and 
\begin{align*}
& \quad \ \tilde\tau_Z\left\{1+ \frac{\tilde\varphi_Z}{n-1} - \tilde\rho^2 \left( 1+ \frac{\tilde\varphi_U}{n-1}\right)^{-1}\right\} 
\\
&= \tilde\tau_Z\left(1+\frac{\tilde\varphi_Z}{n-1} \right) - \tilde\tau_Z\tilde\rho^2\left(1+\frac{b\varphi_U}{n-1}\right)^{-1} 
\\ &= 
\frac{B_2}{1 -\tilde\varphi_Z }\left(1+\frac{\tilde\varphi_Z}{n-1} \right) - \tilde\tau_Z\tilde\rho^2\left(1+\frac{b\varphi_U}{n-1}\right)^{-1}
\\
&= 
B_2 \cdot \frac{(n-1)^{-1}B_2+B_1}{\{(n-1)^{-1} + 1\}B_2}\left\{ 1+\frac{B_1-B_2}{B_2 + (n-1)B_1}\right\} - r\rho^2\tau_Z\left(1+\frac{b\varphi_U}{n-1}\right)^{-1}
\\
&= 
B_2 \cdot \frac{B_2+(n-1)B_1}{nB_2}
\frac{n B_1}{B_2 + (n-1)B_1} - r\rho^2\tau_Z \left(1+\frac{b\varphi_U}{n-1}\right)^{-1}
\\ &= B_1 - r\rho^2\tau_Z \left(1+\frac{b\varphi_U}{n-1}\right)^{-1}
\\
& = \tau_Z\left\{1+\frac{\varphi_Z}{n-1} - \rho^2 \left(1+ \frac{\varphi_U}{n-1}\right)^{-1} \right\}, 
\end{align*} 
which imply that \eqref{eq:fullyconnected-precision-1}
and \eqref{eq:fullyconnected-precision-n-1} hold. 


Finally, we show that there exists $b$ such that the quantities in \eqref{eq:tilde_car_non} are well-defined and $\tilde\beta \ne \beta$. 
Note that $(\tilde\beta, \tilde\varphi_U, \tilde\tau_U, \tilde\varphi_Z, \tilde\tau_Z, \tilde\rho, \tilde\tau_\epsilon)$ are the same as $(\beta, \varphi_U, \tau_U, \varphi_Z, \tau_Z, \rho, \tau_\epsilon)$ when $b=1$. 
Thus, the parameters $(\tilde\beta, \tilde\varphi_U, \tilde\tau_U, \tilde\varphi_Z, \tilde\tau_Z, \tilde\rho, \tilde\tau_\epsilon)$ are well-defined, continuous in $b$ when viewed as functions of $b$, and lead to a valid model, as long as $b$ is in a sufficiently small neighborhood of $1$. 
Note that 
\begin{align*}
    \tilde\beta - \beta & = 
    \rho\sqrt{\frac{\tau_Z}{\tau_U}}(1-\varphi_U)^{-1} - r\rho\sqrt{\frac{\tau_Z}{\tau_U}}(1-b\varphi_U)^{-1}
    \\ &= 
    \rho\sqrt{\frac{\tau_Z}{\tau_U}} 
    \left\{ (1-\varphi_U)^{-1}  - r (1-b\varphi_U)^{-1}\right\}\\
    & = 
    \rho\sqrt{\frac{\tau_Z}{\tau_U}} 
    \left\{ (1-\varphi_U)^{-1}  - \frac{\left(1 + \frac{\varphi_U}{n-1} \right)^{-1} - \left(1-\varphi_U\right)^{-1}}{\left(1 + \frac{b\varphi_U}{n-1} \right)^{-1} - \left(1-b\varphi_U\right)^{-1}}  (1-b\varphi_U)^{-1}\right\}\\
    & = \rho\sqrt{\frac{\tau_Z}{\tau_U}} \cdot
    \frac{
    (1-\varphi_U)^{-1}\left(1 + \frac{b\varphi_U}{n-1} \right)^{-1} - \left(1 + \frac{\varphi_U}{n-1} \right)^{-1} (1-b\varphi_U)^{-1}
    }{\left(1 + \frac{b\varphi_U}{n-1} \right)^{-1} - \left(1-b\varphi_U\right)^{-1}}\\
    & = 
    \rho\sqrt{\frac{\tau_Z}{\tau_U}} \cdot
    \frac{
    (1-\varphi_U)^{-1} \cdot
    \frac{1-b\varphi_U}{ 1 + \frac{b\varphi_U}{n-1} }
     - \left(1 + \frac{\varphi_U}{n-1} \right)^{-1} 
    }{\frac{1-b\varphi_U}{ 1 + \frac{b\varphi_U}{n-1} } - 1}.
\end{align*}
Because $\varphi_U \ne 0$ and $\rho \ne 0$, we can verify that $ \tilde\beta - \beta$ is a strictly monotone function of $\frac{1-b\varphi_U}{ 1 + \frac{b\varphi_U}{n-1} }$, which itself is also a strictly monotone function of $b$. 
Thus, $\tilde\beta - \beta$ is strictly monotone in $b$ for $b$ in a sufficiently small neighborhood around $1$. 
Thus, $\tilde\beta - \beta$ must be nonzero for $b\ne 1$. 


From the above, we derive Proposition \ref{prop:car-fullyconnected-nonid}. 
\end{proof}

\end{document}
